\documentclass[10pt]{article}
\usepackage[margin=20mm]{geometry}
\usepackage{amsmath,amssymb,mathtools,fontspec,unicode-math,xurl,hyperref}
\IfFontExistsTF{DejaVu Serif}{\setmainfont{DejaVu Serif}}{\setmainfont{Latin Modern Roman}}
\setmathfont{latinmodern-math.otf}
\setlength{\parindent}{0pt}\setlength{\parskip}{6pt}\sloppy
\title{Technical verification of
a determinant majorization theorem}
\author{AI-assisted technical review}
\date{}
\begin{document}
\maketitle
Abstract: The main determinant inequality in the acquired first arXiv version is supported by a direct coefficient argument, and its gradient formulation admits a separate verification. Four local findings concern radial derivatives, two assertions about general linear pullbacks, and the normalization of the Lagrangian eigenvalues. Explicit examples meet the relevant stated assumptions. The valid restricted constructions and their independence from the main proof are recorded. A proposed objection about completeness is rejected after checking the authors’ definition. The supported assessment is Minor; scientific adjudication remains pending.\par
\section*{Identity, scope and conventions}
This report concerns F. Reese Harvey and H. Blaine Lawson, Jr., A definitive determinant majorization result for nonlinear operators, arXiv:2407.05408v1 [HL]. The acquired PDF has 13 pages, an arXiv watermark dated 7 July 2024, and a printed date of 9 July 2024. Its title includes determinant although the arXiv landing title omits that word. The frozen journal identity is Duke Mathematical Journal 174 (2025), 2749–2763, DOI https://doi.org/10.1215/00127094-2025-0003. The 15-page journal version was not acquired or certified equivalent. References to pages below mean one-based acquired PDF pages; printed page numbers agree.\par
All 13 PDF pages and the complete released TeX source were read. The review covers Definitions 1.1–1.2, Theorems 1.3 and 2.3, Proposition 1.6 and Corollary 1.7, the coefficient and hyperbolicity lemmas, and all the constructions in Section 3. Three native bounded checks were integrated into the primary analysis. Author TeX was inspected as text and was not executed.\par
The inner product on \(\mathcal S(n)\) is \(\langle A,B\rangle=\operatorname{tr}(AB)\). A Gårding–Dirichlet polynomial \(\mathfrak g\) has \(\mathfrak g(I)>0\), is hyperbolic in direction \(I\), and has a Gårding cone \(\Gamma\) containing every positive definite matrix. The central-ray condition is \(D_I\mathfrak g=kI\) with \(k>0\). These conditions, including the strict positivity of \(k\), are retained in the checks below.\par
\section*{Verification of the main argument}
Equation reference: eq:jensen. For a homogeneous diagonal restriction \(p(x)=\sum_{|\alpha|=N}a_\alpha x^\alpha\) with \(a_\alpha\geq0\), put \(w_\alpha=a_\alpha/p(e)\), where \(e=(1,\ldots,1)\). The central-ray condition and Euler’s identity give \[\sum_\alpha w_\alpha=1,\qquad
 \sum_\alpha w_\alpha\alpha_j=\frac{N}{n}.\] Consequently, for \(x_j>0\), concavity of the logarithm gives \[\log\frac{p(x)}{p(e)}
 \geq \sum_\alpha w_\alpha\log x^\alpha
 =\frac{N}{n}\sum_j\log x_j .
 \] This proves Basic Lemma 2.1 with exactly its stated constant. Equality occurs when \(\alpha\cdot\log x\) is constant on the positive coefficient support. No claim that equality forces a scalar matrix is justified or needed: \(\mathfrak g=\det\) has equality for every \(A>0\). The condition \(k>0\) excludes degree zero; degree one causes no difficulty in this argument.\par
For Lemma 2.5, the endpoint passage can be made precise by holding \(\epsilon>0\) fixed throughout the \(N\) directional derivatives in directions \(e_j+\epsilon e\). Repeated interlacing leaves, after \(N-1\) steps, a nonzero linear polynomial positive on the positive orthant. Its derivative in the last positive direction is positive. Taking the limit of this entire multilinear differential expression as \(\epsilon\downarrow0\) gives the desired nonnegative coefficient. One need not regard its penultimate coefficient vector as independent of \(\epsilon\). This clarifies a valid proof.\par
Orthogonal conjugation preserves the central-ray condition by \(D_I\mathfrak g_h=h^t(D_I\mathfrak g)h\). Every diagonal restriction of \(\mathfrak g_h\) has the nonnegative coefficients just proved. Diagonalizing \(A>0\) and applying [eq:jensen] (\#eq:jensen) proves Theorem 1.3. Its constant is sharp at \(A=I\). Continuity also gives the corresponding nonnegative-matrix boundary inequality, although the stated theorem only quantifies over \(A>0\). The proof uses standard hyperbolicity-cone and interlacing properties; these foundational results are not independently reproved in full here.\par
There is a direct check of Proposition 1.6, so the equivalence need not rest solely on the cited forthcoming reference. Write \(f=\mathfrak g^{1/N}\) and \(H_B=D_Bf\) for \(B\in\Gamma\). Gårding concavity and degree-one homogeneity yield \(f(A)\leq\operatorname{tr}(H_BA)\) for \(A>0\). If \(f(A)\geq c(\det A)^{1/n}\), determinant-one positive matrices force \(H_B>0\): a negative eigenvalue makes the trace unbounded below, and a zero eigenvalue permits traces tending to zero after the remaining eigenvalues of \(A\) tend to zero. The matrix arithmetic–geometric mean identity then gives \[c\leq\inf_{\substack{A>0\\\det A=1}}\operatorname{tr}(H_BA)
       =n(\det H_B)^{1/n}.\] Conversely, \(H_A>0\) and \(\det H_A\geq\gamma\) imply, by Euler’s identity, \(f(A)=\operatorname{tr}(H_AA)\geq n\gamma^{1/n}(\det A)^{1/n}\). Thus the constant in Corollary 1.7 is \(\gamma=\mathfrak g(I)^{n/N}/n^n\), as printed.\par
\section*{Local findings and bounded repairs}
Problem 1 (Radial differentiation requires the missing hypotheses).  Proposition 3.1(2), p. 8, asserts that the radial derivative of a Gårding–Dirichlet polynomial is \(I\)-central, without requiring the input to be \(I\)-central. Its proof on p. 9 explicitly uses that condition.\par
On \(\mathcal S(2)\) let \(B=\operatorname{diag}(1,2)\), \(\ell(A)=\operatorname{tr}(BA)\) and \(\mathfrak g(A)=\ell(A)^2\). Then \(\mathfrak g(I)=9\) and \(\mathfrak g(tI+A)=(3t+\ell(A))^2\) has only real roots. Its cone is \(\{\ell>0\}\) and contains all \(A>0\). The radial derivative is \(q(A)=6\ell(A)\), but \(D_Iq=6B\) is not a scalar multiple of \(I\). The claimed centrality therefore fails for an admissible input. If the conclusion were used to apply the main theorem to \(q\), it would give a false bound: for \(A=\operatorname{diag}(3/2,2/3)\), \[\det A=1,\qquad q(A)=17<18=q(I).\] There is also a degree endpoint. The \(I\)-central polynomial \(\mathfrak g(A)=\operatorname{tr}A\) on \(\mathcal S(2)\) has radial derivative \(2\), whose gradient is zero and hence does not meet Definition 1.2.\par
A precise valid version assumes an \(I\)-central input of degree \(N\geq2\). Set \(q(A)=d\mathfrak g(A+tI)/dt|_{t=0}\). Interlacing proves its Gårding–Dirichlet property, and differentiated Euler homogeneity gives \[D_Iq=(N-1)D_I\mathfrak g=(N-1)kI.\] Repeated radial differentiation must stop before reaching degree zero. The printed proof’s use of centrality and its stop at degree one are counterevidence showing the intended restricted construction. The finding is Minor: it concerns a supplementary closure statement and is absent from the verified main-theorem dependency path.\par
Problem 2 (A pullback does not test hyperbolicity outside its range).  Proposition 3.1(3)(a), p. 11, claims a two-way implication for arbitrary \(L\in\operatorname{End}(\mathcal S(n))\). A non-surjective \(L\) only tests directions in its range.\par
Take \(n=3\), \(\mathfrak g(A)=\det A\), \(J=\operatorname{diag}(1,-1,-1)\), \(L(A)=\frac{\operatorname{tr}A}{3}J\), and \(B=I\). The original \(\mathfrak g\) is itself \(I\)-central and Gårding–Dirichlet. Since \(\det J=1\), \[\mathfrak g_L(A)=\left(\frac{\operatorname{tr}A}{3}\right)^3,\qquad
 \mathfrak g_L(tI+A)=\left(t+\frac{\operatorname{tr}A}{3}\right)^3.\] Thus \(\mathfrak g_L\) is \(I\)-hyperbolic and \(\mathfrak g_L(I)=1\). But, with \(C=E_{12}+E_{21}\), \[\det(tJ+C)=t(t^2+1),\] so \(\mathfrak g\) is not hyperbolic in direction \(L(I)=J\). All hypotheses of the displayed assertion are met; the positivity condition introduced separately in part (e) is not an assumption of part (a).\par
The implication from ambient hyperbolicity to pullback hyperbolicity is valid. The reverse implication becomes valid if \(L\) is surjective, because then every ambient test matrix is \(L(A)\). The eigenvalue and cone identities in (b)–(c) are valid under the ambient hyperbolicity assumption. The invertible congruences used in Lemma 2.4 meet it; hence this Minor error in the general extension does not invalidate that lemma or Theorem 1.3.\par
Problem 3 (Closed-cone preservation can annihilate the polynomial).  Proposition 3.1(3)(e), p. 11, assumes \(L(\mathcal P)\subset\overline{\Gamma(\mathfrak g)}\) and asserts that \(\mathfrak g_L\) is Gårding–Dirichlet. The closure bar was checked in the actual PDF. The proof on p. 12 instead uses strict positive definiteness of \(L(P)\).\par
On \(\mathcal S(2)\) take \(\mathfrak g=\det\) and \(L(A)=a_{11}E_{11}\). For every positive semidefinite \(P\), \(L(P)\) is positive semidefinite, so the stated closed-cone condition holds. Nevertheless \(\mathfrak g_L\equiv0\), and in particular \(\mathfrak g_L(I)=0\), contrary to Definition 1.1.\par
One sufficient repair is to retain closed-cone preservation and add \(L(I)\in\Gamma(\mathfrak g)\). For \(P>0\) choose \(\epsilon>0\) with \(P-\epsilon I\geq0\). Then \[L(P)=\epsilon L(I)+L(P-\epsilon I)
       \in\Gamma(\mathfrak g)+\overline{\Gamma(\mathfrak g)}
       \subset\Gamma(\mathfrak g).\] Hyperbolicity in direction \(L(I)\) and the valid forward pullback argument now give exactly the required conclusion. There is an important alternative reading: the introductory assertion that the pullback has degree \(N\) might be intended to impose \(\mathfrak g_L\not\equiv0\) as an additional assumption. Under that reading the counterexample is excluded. Indeed, closed-cone preservation and \(\mathfrak g_L\not\equiv0\) imply \(L(I)\in\Gamma\): choose \(P_0>0\) with \(\mathfrak g_L(P_0)\neq0\), hence \(L(P_0)\in\Gamma\), and then choose \(c>0\) with \(cI-P_0\geq0\). It follows that \(cL(I)\in\Gamma+\overline\Gamma\subset\Gamma\). Thus the finding is specifically the missing nondegeneracy condition in the stated construction for arbitrary \(L\), not a failure under that stronger interpretation. The later condition \(L(I)=kI\), \(k>0\), in part (f) would supply this missing condition if explicitly imposed on part (e). Moreover, the chain rule and trace-preserving argument of part (f) are valid. These are substantive restrictions on the counterexample’s scope. The assessment is Minor, since the faulty general construction is not used by the verified main proof.\par
Problem 4 (The Lagrangian factors are not normalized \(I\)-eigenvalues).  In the proof of Lemma 3.5, p. 10, equation (3.3) calls \(\Lambda_\pm(A)=\frac12\operatorname{tr}A\pm\lambda_1\pm\cdots\pm\lambda_n\) the \(I\)-eigenvalues on \(\mathcal S(2n)\).\par
For any homogeneous hyperbolic polynomial, \(\mathfrak g(tI+I)=\mathfrak g(I)(t+1)^N\), so every defined \(I\)-eigenvalue of \(I\) equals \(1\). The displayed \(\Lambda_\pm(I)\) instead equals \(n\), since the skew-Hermitian part of \(I\) vanishes. For \(n>1\) these cannot be the normalized \(I\)-eigenvalues.\par
The product factors obey \(\Lambda_\pm(tI+A)=nt+\Lambda_\pm(A)\). Thus the true \(I\)-eigenvalues are \(\Lambda_\pm(A)/n\), and their sum is \(2^{n-1}\operatorname{tr}A/n\). This remains a positive multiple of the ordinary trace, so the central-ray conclusion survives. The separately computed \(\mathfrak g(I)^{1/N}=n\) and the sharp inequality (3.2) are unchanged. The verified discrepancy is a Minor normalization error in the proof, not a failure of the stated Lagrangian determinant inequality.\par
\section*{Counterevidence, dependence and assessment}
The claim in Remark 3.2 about completeness was checked against the authors’ earlier definition [Retained latex: \textbackslash{}cite[Appendix B]\{HLprevious\}]. It concerns all the variables in the underlying \(\mathbb R^n\), equivalently \(\mathcal P\cap E=\{0\}\) for the edge \(E\). It does not require the polynomial to have zero edge in the entire matrix space. The proposed trace-power objection is therefore rejected. Indeed, if \(0\neq P\geq0\) belonged to the edge, then \(D_I\mathfrak g(P)=0\) but centrality gives \(k\operatorname{tr}P>0\). The earlier source explicitly distinguishes this notion from regularity of the cone.\par
The product construction in Proposition 3.1(1) also passes: the hyperbolic roots combine, the cones intersect, and \(D_I(\mathfrak gh)=h(I)D_I\mathfrak g+\mathfrak g(I)D_Ih\) is a positive scalar multiple of \(I\). The nonhyperbolic example \(\|A\|^2\det A\) in Remark 3.3 satisfies the broader coefficient theorem in dimension at least two; it does not contradict the separate Gårding hypotheses.\par
The supported overall assessment is Minor. This rests on an explicit main-proof verification and dependency separation. It does not rest on the repairs being brief. No counterexample to Theorems 1.3 or 2.3, Corollary 1.7, or the sharp Lagrangian inequality has been found within this inspection. No general PDE existence or interior-regularity theorem is independently certified here.\par
Item | Judgment | Repair and remaining scope\par
Radial derivative | Minor | Central input and \(N\geq2\) verified sufficient.\par
Pullback converse | Minor | Forward implication valid; surjectivity restores converse.\par
Closed-cone map | Minor | Add \(L(I)\in\Gamma\); proof supplied.\par
Lagrangian roots | Minor | Divide by \(n\); sharp bound unchanged.\par
Completeness | Rejected | Existing definition and direct edge argument support the claim.\par
Author adoption is unknown and has not been requested. The current arXiv record lists v1; a bounded correction/material search identified no separate supplement or computational release. This is not an exhaustive claim about later corrections. The missing journal text and external PDE developments remain outside the version-specific conclusion.\par
\section*{Evidence and review record}
The accompanying evidence contains the source PDF and manifest, complete coverage and findings records, proof notes and exact calculations, actual source-page inspection receipts, and the report’s current TeX, compiler, dependency, cache, log, PDF and image bindings. Every final report page is actually inspected after compilation. Failed attempts and predecessor records are retained. Rendering alone is not recorded as inspection.\par
The primary executor integrated three native in-conversation bounded checks. These are AI checks, not human review or independent acceptance. The controller’s UI observation and runtime self-description are separately attributed in provenance; backend equivalence is not asserted. Scientific decision is `pending`; `human\_reviews=[]`. No author code was run.\par
2 F. R. Harvey and H. B. Lawson, Jr., A definitive determinant majorization result for nonlinear operators, arXiv:2407.05408v1 (2024), acquired PDF and released source. https://arxiv.org/abs/2407.05408v1. F. R. Harvey and H. B. Lawson, Jr., Determinant majorization and the work of Guo–Phong–Tong and Abja–Olive, Calc. Var. Partial Differential Equations 62 (2023), article 153, DOI https://doi.org/10.1007/s00526-023-02485-8. The completeness definition and Appendix B were consulted; the entire earlier paper was not subjected to a new review.\par
\end{document}
